\section{Формальная верификация и исполнимые свидетели}\label{sec:verification}
Утверждения разделов~\ref{sec:investigation}--\ref{sec:closure} опираются на три различных вида свидетельств: термы доказательств, проверенные ядром Lean, исполнимые примеры конечной модели и независимые эталонные вычисления. Ни одно из них не является измерением физического быстродействия, а согласие конечного числа запусков не заменяет общих утверждений, доказанных в формализации.

\subsection{Механизация и границы свидетельств}
Принятый канонический стек использует Lean~4.36.0 \cite{lean4}, закреплённый за коммитом\\
\texttt{611c117bfaf7b7ae1cfcf246ce6fba31e8f509e9}. Он включает 21 исходный файл Lean (включая CLI) и 542 явно аудируемые декларации: 520 теорем и 22 определения. Таблица~\ref{tab:formal-audit} разбивает цели аудита по группам зависимостей архитектуры; эти числа не являются числом независимых результатов или мерой сложности доказательств. Последние два файла предоставляют типизированный релиз и его исполнимую проекцию. Две ранние исторические ветви сохранены отдельно и не импортируются этим стеком.

\begin{table}[htbp]
\centering\small
\caption{Соответствие артефактов: исторические идентификаторы этапов, математические конструкции и цели аудита. Раздел~2 имеет отдельные исторические доказательные артефакты.}\label{tab:formal-audit}
\begin{tabularx}{\linewidth}{@{}lXrr@{}}\toprule
Этап & Основная конструкция & Файлы & Цели\\\midrule
T19--T22 & Исследование, перенос, библиотека, открытие & 4 & 132\\
T22.9--T24.7 & Платный поиск, отсечение, продолжение, ядра, обучение & 6 & 119\\
T25--T27 & Метаисследование и свидетельства для решения & 3 & 101\\
T28--T29c & Замена контроллера и законы трасс & 4 & 132\\
T30--T31 & Интерпретация политик и замыкание & 2 & 40\\
T32 & Релиз и CLI & 2 & 18\\\midrule
Всего & & 21 & 542\\\bottomrule
\end{tabularx}
\end{table}

Канонический исполнитель проверяет полный манифест и 19 закреплённых унаследованных исходников, компилирует в чистом временном каталоге импортов, аудирует ожидаемые полные имена и зависимости каждой заданной декларации, запускает 13 независимых скриптов конечных моделей и сравнивает фактический вывод CLI с отдельным эталоном Python. Единственные разрешённые аксиоматические зависимости 542 целей --- \rc{propext} и \rc{Quot.sound}; предоставленный успешный журнал фиксирует ноль ошибок и предупреждений компилятора, отсутствие зависимостей от \rc{sorryAx} и \rc{Classical.choice} в этих отчётах. Два анализатора требуют точных имён отчётов и отклоняют неполные отчёты, дубликаты, неразрешённые зависимости и неуспешные коды выхода компилятора. Десять тестовых примеров анализатора аудита и двенадцать примеров анализатора проекций проверяют эти механизмы; это тесты проверяющей программы, а не дополнительные запуски ядра.

Принятый успешный журнал предоставлен автором. В среде первой сборки архива не было Lean, и запуск остановился с кодом~127. Позднее предоставленные полные журналы модулей и исполнителя побайтно совпадают с соответствующим прежним успешным журналом. Это согласуется с детерминированным повтором, но независимо не устанавливает отдельного нового процесса компиляции. Согласование сохраняет это различие в \rc{catalog/VERIFICATION.json} со статусом \rc{SUPPLIED\_LOGS\_AND\_PYTHON\_REPLAY\_PASS}. Последующий повторный аудит согласованного архива (повторённый при подготовке английской редакции RC4) проверил все 21 закреплённый исходник и 542 имени отчётов, 21 журнал модулей и девять проецируемых сценариев релиза; все 45 команд Python завершились с нулевым кодом. Этот аудит не вызывал Lean. Совпадение хешей устанавливает идентичность содержимого, а не независимое удостоверение исполнения. Историческая нумерация ограничена этой картой артефактов. Ранние исходники и первичные связи T03 и T14--T18 сохранены отдельно; для T01, T02 и T04--T13 нет однозначно установленных первичных связей с этапами. Они не являются предпосылками канонического доказательства релиза. См. \rc{GAPS\_AND\_STATUS\_RU.md}.

\clearpage
\subsection{Положительные и отрицательные вычислительные свидетели}
Девять сценариев релиза проверяют составной исполнитель, а не изолированный предикат допуска. Для каждого независимый эталон строит канонический список ответов, проверяет остатки, упорядоченные квитанции поиска и обновлений, события политики и прежние сертификаты. Сохранённая проекция Lean CLI совпадает с эталоном. Результат \emph{impossible} означает исчерпание с пустым каноническим списком в заданном классе морфизмов, а не невозможность для произвольных алгоритмов.

\begin{table}[htbp]
\centering\small
\caption{Свидетели релиза из \rc{FINITE\_RELEASE.json} и предоставленного журнала CLI. Все начинаются с нулевых записанных расходов, кроме случая Predecessor с пятью единицами расходов. $B_0$ --- начальный доступный остаток, $D_f$ --- конечные совокупные расходы, $B_f$ --- конечный остаток. Последний столбец --- число канонических сертификатов.}\label{tab:release-cases}
\begin{tabularx}{\linewidth}{@{}lXrrrr@{}}\toprule
Сценарий & Вердикты метаобновлений / корневой исход & $B_0$ & $D_f$ & $B_f$ & $|A_f|$\\\midrule
Baseline & нет; сохранить & 110 & 7 & 103 & 0\\
Single & принято; заменить & 110 & 6 & 104 & 0\\
Mixed & принято, отклонено; заменить & 110 & 7 & 103 & 0\\
Double & принято, принято; заменить & 110 & 6 & 104 & 0\\
Fallback & отложено, отложено; блокировка & 100 & 7 & 93 & 0\\
Liquidity baseline & нет; заменить & 103 & 4 & 99 & 0\\
Liquidity rejected & отклонено; отказ по ликвидности & 103 & 10 & 93 & 0\\
Nonempty & нет; сохранить & 100 & 11 & 89 & 2\\
Predecessor & нет; сохранить & 95 & 11 & 89 & 2\\\bottomrule
\end{tabularx}
\end{table}

\paragraph{Сквозной пример по слоям.}
Рассмотрим действительный свидетель препятствия по глубине: $S,T$ имеют по одному миру, два порта, одно терминальное действие, постоянные ложные чтения, нулевые терминальные потери и цены портов; $d_S=1$ и $d_T=0$. Тест морфизма раздела~3 не проходит ограничение глубины, поэтому $\CanonRef=[]$. Тем не менее фронтир без отсечения должен обработать три шага разбиения и четыре шага подготовки по единичной цене (стоимость $7$). Сертифицированное правило глубины отвергает начальную область за один оплаченный шаг (стоимость $1$). При $B=110$, $R=R'=100$ и $H=K=1$ условия раздела~6 допускают замену: $101\le110$, $102\le110$ и $1+1+1<7$. Корневая замена с завершением стоит $3$, не меняя требуемый результат.

Теперь сравним альтернативные исполнения на той же паре, а не последовательные стоимости. Исходная метаполитика сохраняет старый контроллер (стоимость $7$). Замкнутая модель приобретения сопоставляет ожидающий/исчерпанный фронтир и действия keep/attempt; выбранное ею дерево зондирует, остались ли задания, и пытается заменить контроллер при истинном ответе. Метаобновление оплачивается, его точные границы равны $7$ и $6$, оно устанавливает это дерево. Фактический релиз с одним обновлением платит по единице за обновление, зондирование, подготовку, проверку, миграцию и окончательный вызов правила: всего $6$, остаток $104$, пустой фронтир и сертифицированная невозможность. Это случай \emph{single} таблицы~\ref{tab:release-cases}; отдельные случаи \emph{nonempty} и \emph{predecessor} проверяют положительное каноническое членство и сохранение прежних сертификатов. См. \rc{ControllerReplacement.Witnesses}, \rc{MetaPolicy.Witnesses} и \rc{Release.Witnesses.costs}.

Случай двойного обновления начинается с политики избыточных зондирований (стоимость $8$), а не с исходной политики; её стоимость уменьшается $8\to7\to6$. Обязательное обучение проигрывает фиксированному контроллеру правила чтения во всех 515 финансируемых случаях стенда обучения. Смешанные корневые трассы завершаются со стоимостями 4 и 8 против 7; стенд замыкания проверяет 965 завершённых запусков и 2418 переходов. Пара по ликвидности ($4$ против $10$) опровергает безусловную монотонность на уровне политик. Это конечные модельные свидетели, а не выборки из прикладной генеральной совокупности или измерения физической производительности.

\newpage
\subsection{Артефакт и воспроизведение}\label{sec:artifact}
Приложение --- \rc{CRTT-v1.0-FULL-ARCHITECTURE-RECONCILED.zip} \cite{crttartifact} с неизменённым принятым каноническим содержимым, историческими исходниками, каталогами, полными предоставленными журналами, результатами повторного аудита и скриптами проверки. Его внешний SHA-256:
\begin{center}\small\texttt{e09580cdb8a8f900401390db394a7bde07bb328449b50b059f2a8cd3c175c79e}\end{center}
Полный манифест содержит 766 записей, охватывающих 767-файловую оболочку за исключением самого списка контрольных сумм. Из корня распакованного архива используйте новый выходной каталог для повторного аудита:
\begin{verbatim}
python3 verify_archive.py
python3 replay_evidence.py --output ../CRTT-log-replay
\end{verbatim}
Для свежей проверки ядром войдите в \rc{canonical/} с закреплённым компилятором Lean и выполните \verb|bash run_verify.sh|, сохраняя журнал вне этого каталога. Только последняя команда выполняет свежий запуск Lean. Приложение к рукописи включает \rc{verify\_fresh.sh} --- оболочку чистой распаковки, регистрирующую входной хеш, инструментарий, журнал исполнения и статус процесса. Она не заменяет канонический исполнитель. При подготовке английской редакции RC4 её локальное исполнение снова прошло предварительную проверку архива, но остановилось с кодом~127 до компиляции Lean. Свежая независимая проверка ядром и публичный постоянный идентификатор артефакта остаются задачами публикации; ни одна здесь не заявлена выполненной. Карты исходников связывают утверждения с декларациями и записанными строками аудита.
